\chapter*{Chapitre 4 — Production de chaleur}
\addcontentsline{toc}{chapter}{Chapitre 4 — Production de chaleur}

% ============================================================
\section*{Énoncés}
% ============================================================

\exercice{Dimensionnement de la puissance d'une chaudière collective}

Une résidence de 30 logements de 55\,m² chacun est située en zone climatique H2b (Lyon, température de base $T_{base} = -8$\,°C, $T_{int} = 19$\,°C). Le coefficient de déperdition moyen est $G = \SI{1{,}0}{\watt\per\square\metre\per\kelvin}$. Les besoins ECS représentent \SI{10}{\kilo\watt} en pointe. Le coefficient de foisonnement est 0,75.

\begin{enumerate}
  \item Calculer les déperditions totales de la résidence.
  \item Calculer la puissance de production nécessaire (chauffage + ECS).
  \item Choisir entre trois chaudières à condensation : modèle X (\SI{60}{\kilo\watt}), modèle Y (\SI{80}{\kilo\watt}) et modèle Z (\SI{100}{\kilo\watt}). Justifier le choix en tenant compte d'une marge de sécurité de 10\,\%.
\end{enumerate}

% ----------------------------------------------------------

\exercice{Comparaison PAC vs chaudière à condensation en rénovation}

Un propriétaire doit remplacer sa vieille chaudière gaz dans une maison de 120\,m² (zone H1, $T_{base} = -12$\,°C). Les déperditions sont estimées à \SI{10}{\kilo\watt}. Les émetteurs sont des radiateurs haute température (80/60\,°C). Il hésite entre :
\begin{itemize}
  \item Solution A : chaudière à condensation gaz, $\eta_{PCS} = 104\,\%$ (référencé PCS), prix gaz \SI{0{,}11}{\euro\per\kilo\watt\hour}
  \item Solution B : PAC air/eau haute température, SCOP = 2,8 (fonctionnement à 60\,°C), prix élec \SI{0{,}22}{\euro\per\kilo\watt\hour}
\end{itemize}
Consommation annuelle estimée : 15\,000 kWh thermiques.

\begin{enumerate}
  \item Calculer le coût annuel de chauffage pour chaque solution.
  \item Sachant que la PAC coûte \SI{5000}{\euro} de plus à l'installation, calculer le temps de retour sur investissement simple.
\end{enumerate}

% ----------------------------------------------------------

\exercice{Analyse des performances d'une PAC en exploitation}

Une PAC air/eau est relevée lors de sa mise en service par les mesures suivantes :
\begin{itemize}
  \item Puissance électrique mesurée : \SI{3{,}8}{\kilo\watt}
  \item Température départ eau : \SI{45}{\celsius}
  \item Température retour eau : \SI{40}{\celsius}
  \item Débit eau : \SI{1200}{\litre\per\hour}
  \item Température extérieure : \SI{2}{\celsius}
\end{itemize}
$\rho_{eau} = \SI{992}{\kilogram\per\cubic\metre}$ à 42\,°C, $c_p = \SI{4179}{\joule\per\kilogram\per\kelvin}$.

\begin{enumerate}
  \item Calculer la puissance thermique produite.
  \item Calculer le COP mesuré.
  \item Le fabricant annonce COP = 3,8 en A2/W45. Analyser l'écart.
\end{enumerate}

% ----------------------------------------------------------

\exercice{Analyse des fumées et combustion}

Un technicien effectue un contrôle réglementaire d'une chaudière à gaz naturel de \SI{80}{\kilo\watt} installée dans une chaufferie collective. L'analyseur de fumées donne les résultats suivants :
\begin{itemize}
  \item Teneur en CO\textsubscript{2} mesurée : $[\text{CO}_2]_{mes} = 8{,}2\,\%$
  \item Teneur en CO : $[\text{CO}] = 42\,\text{ppm}$
  \item Température des fumées : $T_{fum} = \SI{185}{\celsius}$
  \item Température de l'air comburant : $T_{air} = \SI{18}{\celsius}$
  \item PCI du gaz naturel : $\text{PCI} = \SI{10{,}33}{\kilo\watt\hour\per\cubic\metre}$
\end{itemize}

\textit{Rappels :}
\begin{itemize}
  \item Teneur en CO\textsubscript{2} max (stœchiométrie gaz naturel) : $[\text{CO}_2]_{max} = 11{,}7\,\%$
  \item Excès d'air : $e = \dfrac{[\text{CO}_2]_{max}}{[\text{CO}_2]_{mes}} - 1$
  \item Rendement de combustion (méthode de Siegert simplifiée) : $\eta_{comb} = 1 - \dfrac{A + B \times T_{fum}}{[\text{CO}_2]_{mes}\,\%}$ avec $A = 0{,}68$ et $B = 0{,}007$ pour le gaz naturel
  \item Seuil réglementaire CO : $< 100\,\text{ppm}$ dans les fumées sèches
  \item Rendement annuel saisonnier minimal : 91\,\% sur PCI (arrêté du 2 octobre 2009 pour chaudières $> 70\,\text{kW}$)
\end{itemize}

\begin{enumerate}
  \item Calculer l'excès d'air $e$ (en \,\%).
  \item Calculer le rendement de combustion $\eta_{comb}$ avec la formule de Siegert.
  \item Le taux de CO est-il conforme ? Quelles peuvent être les causes d'un taux de CO élevé ?
  \item Le rendement est-il conforme à la réglementation ?
  \item Quelle économie de gaz (en m\textsuperscript{3}/an et en €) peut-on espérer en optimisant la combustion pour atteindre $[\text{CO}_2] = 10\,\%$ ? \\
  (consommation actuelle : $85\,000\,\text{kWh}_{\text{PCI}}/\text{an}$, prix gaz : $0{,}095\,€/\text{kWh}$)
\end{enumerate}

% ----------------------------------------------------------

\exercice{Dimensionnement d'une installation solaire thermique}

Un hôtel de 20 chambres à Marseille (zone H2d, Sud méditerranéen, DJU $\approx$ 700\,°C·j/an, irradiance $\approx$ 1700\,kWh/m²/an) souhaite installer des capteurs solaires plans pour préchauffer l'ECS. Données :
\begin{itemize}
  \item Besoin ECS journalier : $V_{ECS} = 500\,\text{L/j}$ à $T_{puis} = \SI{60}{\celsius}$
  \item Température de l'eau froide : $T_{fr} = \SI{15}{\celsius}$
  \item Irradiance moyenne annuelle sur plan incliné 45°S à Marseille (H2d) : $G_{an} = \SI{1700}{\kilo\watt\hour\per\square\metre\per\year}$
  \item Rendement optique du capteur plan : $\eta_0 = 0{,}79$
  \item Coefficient de pertes : $a_1 = \SI{3{,}8}{\watt\per\square\metre\per\kelvin}$
  \item Température moyenne du fluide capteur : $T_m = \SI{50}{\celsius}$
  \item Température extérieure moyenne annuelle : $T_{ext,moy} = \SI{15}{\celsius}$
  \item Taux de couverture solaire cible : $f = 60\,\%$
  \item Rendement des échangeurs et pertes réseau : $\eta_{sys} = 0{,}85$
  \item Irradiance journalière moyenne à utiliser : $G_{moy} = \SI{5{,}0}{\kilo\watt\hour\per\square\metre\per\day}$
\end{itemize}

\begin{enumerate}
  \item Calculer le besoin thermique journalier pour l'ECS (en kWh/j).
  \item Calculer le besoin thermique annuel total (en kWh/an).
  \item Calculer le rendement annuel moyen du capteur :
  \[
    \eta_{cap} = \eta_0 - \frac{a_1 \times (T_m - T_{ext,moy})}{G_{moy} \times 1000}
  \]
  \item Calculer la surface de capteurs nécessaire pour atteindre le taux de couverture $f = 60\,\%$ :
  \[
    S = \frac{f \times \text{Besoin}_{an}}{G_{an} \times \eta_{cap} \times \eta_{sys}}
  \]
  \item Quel est le volume de ballon de stockage recommandé (règle : 50 à 80\,L/m² de capteur) ?
  \item Comparer le coût évité (énergie solaire produite × prix gaz) à l'investissement de $850\,€/\text{m}^2$ pour calculer le TRI.
\end{enumerate}

\newpage
% ============================================================
\section*{Corrigés}
% ============================================================

\subsection*{Corrigé — Exercice 1}

\textit{Question 1 — Déperditions :}
\[
  S_{totale} = 30 \times 55 = \SI{1650}{\square\metre}
\]
\[
  \dot{Q}_{dep} = G \times S \times (T_{int} - T_{base}) = 1{,}0 \times 1650 \times (19 - (-8)) = 1{,}0 \times 1650 \times 27 = \SI{44\,550}{\watt} \approx \SI{44{,}6}{\kilo\watt}
\]

\[
  \boxed{\dot{Q}_{dep} \approx \SI{44{,}6}{\kilo\watt}}
\]

\textit{Question 2 — Puissance nécessaire :}
\[
  \dot{Q}_{total} = (\dot{Q}_{dep} + \dot{Q}_{ECS}) \times \frac{1}{foisonnement} = (44{,}6 + 10) \times \frac{1}{0{,}75} = 54{,}6 \times 1{,}333 \approx \SI{72{,}8}{\kilo\watt}
\]

\[
  \boxed{\dot{Q}_{total} \approx \SI{72{,}8}{\kilo\watt}}
\]

\textit{Question 3 — Choix :} Avec une marge de sécurité de 10\,\%, la puissance minimale à installer est $72{,}8 \times 1{,}10 \approx \SI{80{,}1}{\kilo\watt}$. Le modèle X (\SI{60}{\kilo\watt}) est insuffisant. Le modèle Y (\SI{80}{\kilo\watt}) est légèrement insuffisant (80 < 80,1). Le \textbf{modèle Z (\SI{100}{\kilo\watt})} est retenu pour offrir une marge opérationnelle confortable (25\,\%) permettant d'absorber un pic ECS simultané sans dégradation.

% ----------------------------------------------------------

\subsection*{Corrigé — Exercice 2}

\textit{Question 1 — Coûts annuels :}

\textit{Solution A (chaudière) :} L'énergie achetée en gaz pour produire 15\,000\,kWh thermiques :
\[
  E_{gaz} = \frac{15\,000}{1{,}04} = 14\,423\,\text{kWh}_{\text{PCS}}
\]
\[
  C_A = 14\,423 \times 0{,}11 = \SI{1587}{\euro\per\an}
\]

\textit{Solution B (PAC) :} L'énergie électrique consommée :
\[
  E_{élec} = \frac{15\,000}{2{,}8} = 5\,357\,\text{kWh}_{\text{élec}}
\]
\[
  C_B = 5357 \times 0{,}22 = \SI{1179}{\euro\per\an}
\]

\[
  \boxed{C_A = \SI{1587}{\euro\per\an} \quad \text{vs} \quad C_B = \SI{1179}{\euro\per\an}}
\]

\textit{Question 2 — Temps de retour :}
\[
  \text{Économie annuelle} = 1587 - 1179 = \SI{408}{\euro\per\an}
\]
\[
  TRI = \frac{5000}{408} \approx \mathbf{12{,}3\,\text{ans}}
\]

\[
  \boxed{TRI \approx 12{,}3 \text{ ans}}
\]

\begin{attention}
Ce temps de retour est à comparer avec les aides disponibles (MaPrimeRénov', CEE) qui peuvent le réduire significativement.
\end{attention}

% ----------------------------------------------------------

\subsection*{Corrigé — Exercice 3}

\textit{Question 1 — Puissance thermique :}
\[
  \dot{m} = \frac{1200 \times 992}{3600 \times 1000} = \SI{0{,}331}{\kilogram\per\second}
\]
\[
  \dot{Q} = \dot{m} \cdot c_p \cdot \Delta T = 0{,}331 \times 4179 \times 5 = \SI{6\,916}{\watt} \approx \SI{6{,}9}{\kilo\watt}
\]

\[
  \boxed{\dot{Q} \approx \SI{6{,}9}{\kilo\watt}}
\]

\textit{Question 2 — COP mesuré :}
\[
  COP_{mes} = \frac{6{,}9}{3{,}8} = \mathbf{1{,}82}
\]

\[
  \boxed{COP_{mes} = 1{,}82}
\]

\textit{Question 3 — Analyse :} Le COP mesuré (1,82) est très inférieur au COP annoncé (3,8 en A2/W45). Pistes de diagnostic :
\begin{itemize}
  \item Vérifier l'état du circuit frigorifique (charge en fluide frigorigène)
  \item Vérifier l'encrassement ou le givrage de l'échangeur air/fluide
  \item Vérifier que la PAC n'est pas en mode dégivrage lors de la mesure
  \item Vérifier le débit d'air sur l'évaporateur (filtre encrassé ?)
  \item Vérifier l'absence de court-circuit d'air chaud recyclé vers l'aspiration
\end{itemize}

% ----------------------------------------------------------

\subsection*{Corrigé — Exercice 4}

\textit{Question 1 — Excès d'air :}
\[
  e = \frac{[\text{CO}_2]_{max}}{[\text{CO}_2]_{mes}} - 1 = \frac{11{,}7}{8{,}2} - 1 = 1{,}427 - 1 = 0{,}427
\]
\[
  \boxed{e = 42{,}7\,\%}
\]
Un excès d'air de 42,7\,\% signifie que la chaudière aspire 42,7\,\% d'air de plus que le strict minimum stœchiométrique. Cet air excédentaire est chauffé inutilement par la combustion puis évacué dans les fumées, ce qui pénalise le rendement.

\textit{Question 2 — Rendement de combustion (Siegert) :}

$T_{fum} = 185\,°C$ mais la formule de Siegert utilise $\Delta T = T_{fum} - T_{air} = 185 - 18 = 167\,\text{K}$ intégré dans les coefficients $A$ et $B$ :
\[
  \eta_{comb} = 1 - \frac{A + B \times \Delta T}{[\text{CO}_2]_{mes}\,\%} = 1 - \frac{0{,}68 + 0{,}007 \times 167}{8{,}2}
\]
\[
  = 1 - \frac{0{,}68 + 1{,}169}{8{,}2} = 1 - \frac{1{,}849}{8{,}2} = 1 - 0{,}2255 = 0{,}7745
\]
\[
  \boxed{\eta_{comb} \approx 77{,}5\,\%}
\]

\begin{attention}
Ce rendement de combustion (77,5\,\%) est anormalement bas pour une chaudière récente. Il traduit une perte sensible dans les fumées trop chaudes combinée à un excès d'air excessif. Un réglage fin du brûleur est indispensable.
\end{attention}

\textit{Question 3 — Conformité CO :}
\[
  [\text{CO}] = 42\,\text{ppm} < 100\,\text{ppm} \quad \Rightarrow \quad \textbf{Conforme}
\]
Le taux de CO est réglementairement acceptable. Néanmoins un CO élevé peut provenir de :
\begin{itemize}
  \item Manque d'air (richesse excessive) — inverse du cas présent
  \item Mauvaise atomisation du combustible / injecteur encrassé
  \item Flamme froide (mauvais réglage, encrassement chambre de combustion)
  \item Température de chambre insuffisante (court-circuit de flamme sur paroi froide)
  \item Défaut d'étanchéité du circuit de fumées (faux air diluteur)
\end{itemize}

\textit{Question 4 — Conformité réglementaire du rendement :}

L'arrêté du 2 octobre 2009 impose un rendement saisonnier $\geq 91\,\%$ sur PCI pour les chaudières $> 70\,\text{kW}$.
\[
  \eta_{comb} = 77{,}5\,\% < 91\,\% \quad \Rightarrow \quad \textbf{Non conforme}
\]
La chaudière doit être réglée (réduction de l'excès d'air, abaissement de la température de fumées) avant de pouvoir être mise en service ou maintenue en exploitation.

\textit{Question 5 — Économie potentielle :}

En portant $[\text{CO}_2]$ de 8,2\,\% à 10\,\%, on recalcule le rendement cible :
\[
  e_{cible} = \frac{11{,}7}{10} - 1 = 0{,}17 \quad \Rightarrow \quad 17\,\% \text{ d'excès d'air}
\]
En supposant que la température de fumées diminue à environ $T_{fum,cible} \approx 160\,°C$ ($\Delta T = 142\,\text{K}$) :
\[
  \eta_{comb,cible} = 1 - \frac{0{,}68 + 0{,}007 \times 142}{10} = 1 - \frac{0{,}68 + 0{,}994}{10} = 1 - \frac{1{,}674}{10} = 83{,}3\,\%
\]
Gain de rendement : $\Delta\eta = 83{,}3 - 77{,}5 = 5{,}8\,\%$

Économie en énergie :
\[
  \Delta E = 85\,000 \times \frac{5{,}8}{77{,}5} \approx 6\,355\,\text{kWh}_{\text{PCI}}/\text{an}
\]

Économie financière :
\[
  \Delta C = 6\,355 \times 0{,}095 = \boxed{\SI{604}{\euro\per\year}}
\]

En volume de gaz ($\text{PCI} = 10{,}33\,\text{kWh/m}^3$) :
\[
  \Delta V_{gaz} = \frac{6\,355}{10{,}33} \approx \boxed{615\,\text{m}^3/\text{an}}
\]

% ----------------------------------------------------------

\subsection*{Corrigé — Exercice 5}

\textit{Question 1 — Besoin thermique journalier ECS :}
\[
  Q_j = \frac{V_{ECS} \times \rho_{eau} \times c_p \times (T_{puis} - T_{fr})}{3\,600}
\]
\[
  = \frac{500 \times 1{,}0 \times 4{,}186 \times (60 - 15)}{3\,600} = \frac{500 \times 4{,}186 \times 45}{3\,600}
\]
\[
  = \frac{94\,185}{3\,600} \approx 26{,}2\,\text{kWh/j}
\]
\[
  \boxed{Q_j \approx 26{,}2\,\text{kWh/j}}
\]
(On utilise $\rho_{eau} = 1\,\text{kg/L}$, $c_p = 4{,}186\,\text{kJ/kg·K}$, le facteur $1/3\,600$ convertit kJ en kWh.)

\textit{Question 2 — Besoin thermique annuel :}
\[
  Q_{an} = Q_j \times 365 = 26{,}2 \times 365 = \boxed{\SI{9\,563}{\kilo\watt\hour\per\year}}
\]

\textit{Question 3 — Rendement annuel moyen du capteur :}

$G_{moy} = 5{,}0\,\text{kWh/m}^2/\text{j} = 5{,}0 \times 1\,000 / 3{,}6\,\text{W/m}^2 \approx 1\,389\,\text{W/m}^2$ (en crête). La formule donnée utilise $G_{moy}$ en kWh/m²/j — on applique la conversion pour obtenir une puissance en W/m² :
\[
  G_{moy,\,W} = \frac{G_{moy} \times 1\,000}{3{,}6} = \frac{5{,}0 \times 1\,000}{3{,}6} \approx 1\,389\,\text{W/m}^2
\]
\[
  \eta_{cap} = \eta_0 - \frac{a_1 \times (T_m - T_{ext,moy})}{G_{moy,\,W}} = 0{,}79 - \frac{3{,}8 \times (50 - 15)}{1\,389}
\]
\[
  = 0{,}79 - \frac{3{,}8 \times 35}{1\,389} = 0{,}79 - \frac{133}{1\,389} = 0{,}79 - 0{,}096 = 0{,}694
\]
\[
  \boxed{\eta_{cap} \approx 69{,}4\,\%}
\]

\textit{Question 4 — Surface de capteurs :}
\[
  S = \frac{f \times Q_{an}}{G_{an} \times \eta_{cap} \times \eta_{sys}} = \frac{0{,}60 \times 9\,563}{1\,700 \times 0{,}694 \times 0{,}85}
\]
\[
  = \frac{5\,738}{1\,002{,}8} \approx \boxed{5{,}7\,\text{m}^2}
\]
En pratique, on retiendra \textbf{6 m²} (en unités de capteurs standard de 2 m² chacun, soit 3 capteurs).

\textit{Question 5 — Volume de ballon de stockage :}

Règle empirique : 50 à 80\,L par m² de capteur :
\[
  V_{ballon} = 50 \text{ à } 80\,\text{L/m}^2 \times 5{,}7\,\text{m}^2 = 285 \text{ à } 456\,\text{L}
\]
\[
  \boxed{V_{ballon} \approx 300 \text{ à } 450\,\text{L}}
\]
On retiendra un ballon de \textbf{400\,L} (volume standard disponible) pour offrir une autonomie suffisante en cas de faible ensoleillement.

\textit{Question 6 — TRI de l'installation :}

Énergie solaire produite annuellement :
\[
  E_{sol} = f \times Q_{an} = 0{,}60 \times 9\,563 = \SI{5\,738}{\kilo\watt\hour\per\year}
\]

Coût évité (énergie qui aurait été fournie par gaz, $\eta_{chaud} = 0{,}90$) :
\[
  \Delta C = \frac{E_{sol}}{\eta_{chaud}} \times \text{prix gaz} = \frac{5\,738}{0{,}90} \times 0{,}095 \approx 6\,376 \times 0{,}095 = \boxed{\SI{606}{\euro\per\year}}
\]

Investissement total :
\[
  I = S \times 850 = 5{,}7 \times 850 = \SI{4\,845}{\euro}
\]
(On retient $S = 5{,}7\,\text{m}^2$ pour le calcul ; en pratique l'investissement sera basé sur la surface réellement posée.)

\[
  TRI = \frac{I}{\Delta C} = \frac{4\,845}{606} \approx \boxed{8{,}0\,\text{ans}}
\]

\begin{methode}
Ce TRI de 8 ans est typique pour une installation solaire thermique en zone méditerranéenne (H2d). Les aides MaPrimeRénov' Copropriété ou les CEE peuvent réduire l'investissement de 30 à 50\,\%, ramenant le TRI à 4--5 ans. En zone H1 (Paris), l'irradiance plus faible allonge le TRI à 10--12 ans sans aides.
\end{methode}
